% SOURCE_AUTHORITY: sga5_fr_workpass.tex, lines 13241--13269 (LF-normalized unit).
% SOURCE_SHA256: 791F4EFFC5E02832D5D77ED03518C8156D6F07E4C8238B03545DB93D883FBB28
\section*{Exposição XII}

\begin{center}
\Large\bf FÓRMULAS DE NIELSEN--WECKEN E DE LEFSCHETZ EM GEOMETRIA ALGÉBRICA
\end{center}

\begin{center}
por A. Grothendieck

redigido por I. BUCUR
\end{center}

\bigskip

\section*{1.}
Sejam $X'$ um esquema próprio sobre o corpo algebricamente fechado $k$ e $G$ um grupo finito que opera sobre $X'$.

Para cada elemento $g \in G$, denota-se por $g_{X'} : X' \to X'$ o automorfismo de $X'$ induzido por $g$.

Suponha que exista um aberto $U'$ de $X'$ estável sob $G$ e tal que $G$ opera livremente sobre $U'$, e seja $\La_n = \bZ/\ell^n\bZ$, onde $\ell$ é um número primo distinto da característica de $k$.

Denotemos por $\La_{n,U',X'}$ o feixe $i_!(\La_{n,U'})$ sobre $X'$, onde $\La_{n,U'}$ é o feixe constante sobre $U'$ associado ao anel $\La_n$ e $i : U' \hookrightarrow X'$ é a imersão canônica de $U'$ em $X'$. O feixe $\La_{n,U',X'}$ é naturalmente munido de uma estrutura de feixe com grupo $G$ de operadores. Por conseguinte, o complexo $H_{n,U',X'} = \R\Gamma_{X'}(\La_{n,U',X'})$ é um complexo de $\La_n[G]$-módulos e, mais precisamente, um complexo perfeito (VIII, 7) de $\La_n[G]$-módulos (IX, 2).

Veremos, a seguir, que, para nosso propósito, é mais conveniente considerar, em vez do complexo $\R\Gamma_{X'}(\La_{n,U',X'})$, seu dual $H_{n,U',X'}^\vee = \R\Gamma_{X'}(\La_{n,U',X'})^\vee$ relativamente ao anel $\La_n$, no sentido das categorias derivadas:
\[
H_{n,U',X'}^\vee = \R\Gamma_{X'}(\La_{n,U',X'})^\vee = \R\Hom_{\La_n}(\R\Gamma_{X'}(\La_{n,U',X'}), \La_n) \quad .
\]

O complexo $\R\Gamma_{X'}(\La_{n,U',X'})^\vee$ é um complexo perfeito de $\La_n[G]$-módulos.
